\section{Conclusion}

This paper provides a comparison of five MILP solvers in the context of cryptanalytic applications, specifically related-key differential attacks on the ITUbee lightweight block cipher. Our findings show that solver selection has a large effect on the feasibility of MILP-based related-key differential cryptanalysis.

Key conclusions include:
\begin{enumerate}
    \item Commercial MILP solvers (Gurobi and CPLEX) are more efficient than open-source alternatives and can be preferred for larger, time-sensitive analyses.
    \item For organizations without access to commercial solvers, SCIP is a good open-source option that performs well on moderate problem sizes.
    \item The performance advantage of commercial solvers grows with problem complexity, suggesting that their methods are more effective on harder instances.
    \item Cryptanalysts should evaluate solver options based on their constraints (computational budget, licensing availability, problem size) when conducting related-key differential cryptanalysis.
\end{enumerate}


